\chapter{Food Allergy Testing}
\section*{What Is This Test?} Food allergy testing evaluates whether immune reactions to specific foods could explain reproducible symptoms.
\section*{Alternative Names} Food allergy blood test; food skin-prick test; oral food challenge.
\section*{Principle} Skin-prick tests and serum specific-IgE assays detect sensitization; a supervised oral food challenge determines clinical reactivity when appropriate.
\section*{Why This Test Is Ordered} It is used for immediate hives, swelling, wheezing, vomiting, or anaphylaxis plausibly linked to a food.
\section*{Primary vs Secondary Test} History is primary. Tests support the evaluation; broad panels and IgG food tests do not diagnose food allergy.
\section*{How This Test Is Performed} Selected foods are tested on skin or in blood. An oral challenge gives measured food portions under medical supervision with emergency treatment available.
\section*{What Preparation Is Needed} Antihistamines may affect skin testing and should be managed only as directed. Never perform a food challenge at home when allergy is suspected.
\section*{Understanding Results} A positive skin or blood test indicates sensitization, not necessarily clinical allergy. A negative result lowers risk but must fit the history.
\section*{Follow-up Tests} Supervised oral challenge, elimination and reintroduction plan, component testing, or allergy referral may follow.
\section*{References} \begin{enumerate}\item MedlinePlus. Food Allergy Testing.\item National Institute of Allergy and Infectious Diseases. Food allergy diagnosis guidance.\end{enumerate}
\clearpage\section*{Understanding Results and Follow-up}
The laboratory report should identify the specimen, method, units, reference interval or decision limit, and whether the result is positive, negative, abnormal, or inconclusive. Reference intervals vary with age, sex, pregnancy, specimen type, assay, and laboratory; a result outside a range is not automatically a diagnosis, and a result within a range does not always exclude disease. Discuss unexpected results with the ordering clinician, who can decide whether repeat or confirmatory testing, treatment, or referral is appropriate. Do not change medicines or delay urgent care based on an isolated result.
\section*{Regulatory Status and Authorization Date}This chapter describes a test category, not a specific commercial product. FDA, CDSCO, EU IVDR, and MHRA approval or registration dates therefore cannot be assigned without the assay manufacturer, product name, instrument, and intended use. For laboratory-developed tests, an individual agency approval date may not exist. See the Preface for the interpretation of regulatory status.
\clearpage
